\chapter{Povzetek dela}

\section{Priprava podatkov}
Trenutno uporabljamo naslednje datasete:\\
\begin{tabular}{|p{9cm}|p{6cm}|}
\hline
\textbf{Ime} & \textbf{Pacienti} \\
\hline
CAP Sleep Database & Odrasli pacienti (Kontrole, Bruksizem, Nespečnost, Narkolepsija, Nočna frontalna lobarna epilepsija, Periodični gibi udov, Vedenjska motnja v spanju REM, Motnje dihanja v spanju) \\
\hline
Dataset of EEG recordings containing HFO markings for 30 pediatric patients with epilepsy & 30 otrok in mladostnikov s fokalno ali generalizirano epilepsijo \\
\hline
NCH Sleep DataBank (Pediatric: newborns to adolescents) & Otroci \\
\hline
Normative longitudinal EEG recordings during sleep stage II in the first year of age & Dojenčki, 15 dni – 12.5 mesecev starosti \\
\hline
Sleep-EDF Database Expanded & Odrasli zdravi pacienti, ne jemljejo zdravil \\
\hline
Haaglanden Medisch Centrum sleep staging database & Odrasli pacienti \\
\hline
Adult High-Density \& Wearable EEG & Odrasli pacienti \\
\hline
\end{tabular}

Na vsakem zapisu izberemo naslednje kanale: F4-C4, C4-P4, P4-O2, F3-C3, C3-P3, P3-O1.
Če kanala ni včasih uporabimo substitucije (trenutno samo na datasetu CAP), sicer se uporablja same ničle za označevanje manjkajočega podatka.

Na vsakem zapisu uporabimo filter ter spremenimo frekvenco vzorčenja na 128 Hz.
\begin{verbatim}
    raw.filter(l_freq=0.3, h_freq=30)
    raw.resample(sfreq=TARGET_SFREQ)
\end{verbatim}

Podatke razdelimo na 30 sekundne intervale ter za vsakega shranimo spektrogram z naslednjimi nastavitvami.

\begin{verbatim}
def get_spectrogram(raw):
        segment_data = raw.get_data()

        spectrograms = []
        for channel_data in segment_data:
            f, t, Zxx = scipy.signal.stft(
                channel_data,
                fs=TARGET_SFREQ,
                nperseg=2 * TARGET_SFREQ,
                noverlap=TARGET_SFREQ
            )
            
            f_mask = (f < 50)
            Zxx = Zxx[f_mask, :]

            spectrogram = np.log1p(np.abs(Zxx))
            spectrograms.append(spectrogram)
        return np.stack(spectrograms, axis=0).astype(np.float32)
\end{verbatim}

    
Funkcija najprej izračuna kratkočasovno Fourierovo transformacijo 
(STFT) za vsak kanal vhodnega signala, kjer za segmentacijo uporabi 
okno dolžine 256 vzorcev (2 * 128 Hz) s 50-odstotnim prekrivanjem 
(128 vzorcev). Nato dobljeni spektrogram omeji le na frekvence pod 
50 Hz, na koncu pa nad absolutnimi vrednostmi izvede še 
logaritemsko transformacijo `log(1+x)`.


\section{Model}
Za klasifikacijo EEG signalov sem uporabil model, 
ki združuje konvolucijsko mrežo (CNN) in Transformer.
CNN deluje kot začetni vmesnik, ki iz spektrogramov
izvleče uporabne značilnice in jih preoblikuje v vektorje.
Ti vektorji nato služijo kot vhod za Transformer encoder.
Pri učenju sem uporabil Focal Loss, ker bolje obravnava 
neuravnotežene razrede, ter Cosine Annealing 
scheduler za prilagajanje hitrosti učenja. 
Podatke sem pripravil v obliki drsečih oken, 
pri čemer oznako določa sredinski del okna, 
poleg tega pa se pri treniranju kanali EEG naključno
premešajo, da model postane bolj robusten. Za stabilnost
sem dodal dropout, validacijo pa spremljal z zgodnjim
ustavljanjem. Na koncu sem rezultate ovrednotil z matriko zmede,
Cohenovo kapo in makro F1 mero.

\section{Eksperimenti in rezultati}
\subsection{Baseline}
Najprej kot vhod preprosto dodamo vse podatke.
Najboljši rezultati:
\begin{verbatim}
Confusion Matrix:
[[ 4547   139   488    74   116]
 [  565   604  1017    20   542]
 [  184   174 11352   810   663]
 [   21     3  1187  4940    21]
 [   56    91   710    35  4806]]
Accuracy: 0.7914669078848183
Cohen's Kappa: 0.714181668325527
F1 Score (macro): 0.7224575201952812
\end{verbatim}

\subsection{Dodajanje podatkov o starosti in bolezni}
Dodamo naslednje embeddinge:

\begin{itemize}
    \item Starost: Dojenček, Otrok, Odrasli
    \item Bolezen: Zdravi, Epilepsija
\end{itemize}
Podatek v model vstavimo šele pred klasifikacijo (glede na teste je delovalo bolje, kot če damo v transformer)
\begin{verbatim}
    age_embedding = self.age_emb(age_group)
    health_embedding = self.health_emb(health_group)

    combined_features = torch.cat([pooled, age_embedding, health_embedding], dim=1)

    # 5. Classify
    logits = self.classification_head(combined_features)
\end{verbatim}
Najboljši rezultati:
\begin{verbatim}
Confusion Matrix:
[[ 4324   311   411    94   224]
 [  446   823   950    23   506]
 [  137   226 10925  1207   688]
 [   15     6   978  5171     2]
 [   23   146   727    58  4744]]
Accuracy: 0.7835670134177597
Cohen's Kappa: 0.705819791369912
F1 Score (macro): 0.7283874026980337
\end{verbatim}

\subsection{Ločeni modeli}
Uporabimo isti model kot za Baseline, treniramo posebej 
za otroke in odrasle. TODO: Razbiti na več poddatasetov.
Za vsak model prilagodimo parametre.

Najboljši rezultati:
\begin{verbatim}
Confusion Matrix:
[[ 4397   268   439   101   159]
 [  489   741   995    38   485]
 [  152   188 10989  1323   531]
 [   20     6   922  5222     2]
 [   50   115   775   100  4658]]
Accuracy: 0.7841700587969245
Cohen's Kappa: 0.7060608603276319
F1 Score (macro): 0.724619531001996
\end{verbatim}